The Schwartzberg score was introduced by \citet{schwartzberg1966reapportionment}
as a measure of district compactness derived from the isoperimetric ratio.
Unlike PP, which maps to $[0,1]$ with $1$ representing a circle,
S maps to $[1,\infty)$ with $1$ representing a circle:
a district with S $= 2$ is ``twice as non-circular'' as a circle
in the sense that its perimeter is twice the perimeter of a circle
with the same area.

\subsection{Why S Is Preferred in Legislative Contexts}

Practitioners and legislators often find S more intuitive than PP because:
\begin{enumerate}
  \item The $[1,\infty)$ scale avoids the counterintuitive PP property that
    ``higher is more compact.'' S $= 1$ is compact; S $= 5$ is very non-compact.
  \item S has the dimension of a ``stretch factor'': S $= 2.0$ means the district's
    boundary is twice as long as a circle of the same area would require.
  \item For statutory compliance, legislators can write ``districts shall have
    Schwartzberg score $\leq 2.0$'' rather than ``Polsby-Popper score $\geq 0.25$,''
    which is more natural in plain English.
\end{enumerate}

\subsection{Colorado's Mandate}

Colorado's constitution (Art.\ V, \S 44, enacted by Amendment Y in 2018)
requires the Independent Congressional Redistricting Commission to consider
compactness as one of its primary criteria.
The commission's 2021 procedural rules operationalised ``compactness'' using
a Schwartzberg-equivalent perimeter-to-area ratio~\citep{co_irc2021}.
This makes Colorado the most clearly documented state in which S-equivalent
notation is the legally operative compactness metric.

\subsection{Algebraic Equivalence}

Despite differences in scale and legislative preference, S and PP carry
identical geometric information.
Section~3 proves $S = 1/\sqrt{\text{PP}}$ exactly, so the \textsc{bisect}
PP output can always be converted to S using this formula.
The conversion table in Section~4 makes this practical.
